\section{The Bard target}

Reassessing Bard et al.\ means first stating, from the article rather than its
reputation, what they actually claimed. The abstract makes four claims for
magnitude estimation: that it can \enquote{solve the measurement scale problems
which plague conventional techniques}, \enquote{provide data which make fine
distinctions robustly enough to yield statistically significant results of
linguistic interest}, be \enquote{usable in a consistent way by linguistically
naive speaker-hearers}, and \enquote{allow replication across groups of
subjects} \citep[32]{BardEtAl1996}.

The core of the case is an argument about scale type, not about reliability.
The trouble with conventional annotation, they argue, is \enquote{the use of the
wrong kind of measurement scale} \citep[32]{BardEtAl1996}: the asterisk
system, like any fixed-point scale, \enquote{predetermines the number of
distinctions subjects may use} and so censors finer ones
\citep[35]{BardEtAl1996}. Magnitude estimation is the proposed remedy
because it \enquote{does not restrict the number of values}, giving subjects
\enquote{complete freedom about which of the infinite set of numbers to use} and
yielding interval and ratio data \citep[41]{BardEtAl1996}.

This sharpens Bard et al.'s claim and its boundary. They do make a
method-superiority claim, but a specific one: an open, ratio-yielding response
recovers gradient distinctions that fixed-point scales lose by construction.
The claim concerns resolution and scale-boundary censoring, and it can be
tested. The broader reading sometimes met in reception, that magnitude
estimation is the privileged formal method in general, goes beyond the article;
Bard et al.\ compare magnitude estimation with informal annotation and
fixed-point scales, not with the later inventory of Likert, forced-choice,
yes/no, and Thurstonian tasks.

Bard et al.\ are also more tentative than the reception. They grant that the
advantages \enquote{are garnered at considerable empirical cost}, and they leave
open which psychological model underlies acceptability, since \enquote{the data
currently do not decide between them} \citep[65]{BardEtAl1996}. The fair
target for reassessment is the resolution claim, not a gold-standard claim the
authors didn't make.

The distinction matters for design. If the live claim is that an open ratio
scale recovers distinctions fixed scales censor, then the later method
comparisons have to test resolution and scale-boundary effects alongside
convergence and reliability. Because Bard participant-level rows aren't
available to this project, the paper can't replicate the original experiment
directly; the reassessment is a conceptual replication and secondary-data test
of whether the resolution claim retains independent evidential force.
